\section{Implementation results}
\label{sec:results}

% Run numbers below come from the run records through
% scripts/report/a2_results_tables.py; regenerate instead of editing them.
% Generated by scripts/report/a2_results_tables.py from the U-Net training summary and the ablation results. Do not edit.
\newcommand{\unetBestEpoch}{38}
\newcommand{\unetBestValMiou}{36.9}
\newcommand{\unetValMiou}{36.94}
\newcommand{\unetValDice}{48.05}
\newcommand{\unetWarmupValMiou}{18.3}
\newcommand{\unetLastGap}{0.04}
\newcommand{\unetFirstTrainLoss}{1.95}
\newcommand{\unetLastTrainLoss}{0.74}
\newcommand{\unetLowestValLossEpoch}{38}
\newcommand{\unetValBelowTrain}{30}
\newcommand{\unetLossGap}{0.02}
\newcommand{\unetZeroClasses}{PF and RO}
\newcommand{\unetTrainMinutes}{14.9}
\newcommand{\unetEpochSeconds}{22}
\newcommand{\unetLatency}{4.8}
\newcommand{\unetInferenceBatch}{4}
\newcommand{\unetPeakMemory}{0.97}
\newcommand{\unetGflops}{28.3}
\newcommand{\deeplabGflops}{95.8}
\newcommand{\flopRatio}{3.4}
\newcommand{\deeplabMillions}{39.64}
\newcommand{\diceTimeOverhead}{1}
\newcommand{\deeplabCeBestEpoch}{31}
\newcommand{\deeplabCeValMiou}{62.05}
\newcommand{\deeplabCeValDice}{73.91}
\newcommand{\deeplabCeTrainMinutes}{21.4}
\newcommand{\deeplabCeLatency}{7.8}
\newcommand{\deeplabCePeakMemory}{1.93}
\newcommand{\deeplabDiceBestEpoch}{24}
\newcommand{\deeplabDiceValMiou}{62.28}
\newcommand{\deeplabDiceValDice}{74.42}
\newcommand{\deeplabDiceTrainMinutes}{21.6}
\newcommand{\deeplabDiceLatency}{7.8}
\newcommand{\deeplabDicePeakMemory}{1.93}
\newcommand{\deeplabGain}{25.11}
\newcommand{\unetMillions}{7.76}
\newcommand{\unetBoundaryIoU}{18.04}
\newcommand{\unetInteriorIoU}{37.88}
\newcommand{\diceDeltaMiou}{0.23}
\newcommand{\diceDeltaDice}{0.51}
\newcommand{\unetIouBW}{80.62}
\newcommand{\deeplabCeIouBW}{83.89}
\newcommand{\deeplabDiceIouBW}{83.84}
\newcommand{\diceDeltaBW}{0.05}
\newcommand{\unetIouHD}{30.83}
\newcommand{\deeplabCeIouHD}{70.32}
\newcommand{\deeplabDiceIouHD}{70.17}
\newcommand{\diceDeltaHD}{0.15}
\newcommand{\unetIouPF}{0.00}
\newcommand{\deeplabCeIouPF}{11.94}
\newcommand{\deeplabDiceIouPF}{14.87}
\newcommand{\diceDeltaPF}{2.93}
\newcommand{\unetIouWR}{41.16}
\newcommand{\deeplabCeIouWR}{66.88}
\newcommand{\deeplabDiceIouWR}{66.42}
\newcommand{\diceDeltaWR}{0.45}
\newcommand{\unetIouRO}{0.00}
\newcommand{\deeplabCeIouRO}{70.58}
\newcommand{\deeplabDiceIouRO}{68.20}
\newcommand{\diceDeltaRO}{2.38}
\newcommand{\unetIouRI}{67.94}
\newcommand{\deeplabCeIouRI}{75.36}
\newcommand{\deeplabDiceIouRI}{75.28}
\newcommand{\diceDeltaRI}{0.07}
\newcommand{\unetIouFV}{33.14}
\newcommand{\deeplabCeIouFV}{66.89}
\newcommand{\deeplabDiceIouFV}{67.13}
\newcommand{\diceDeltaFV}{0.23}
\newcommand{\unetIouSR}{41.85}
\newcommand{\deeplabCeIouSR}{50.56}
\newcommand{\deeplabDiceIouSR}{52.30}
\newcommand{\diceDeltaSR}{1.74}
\newcommand{\ablationFrozenDelta}{7.65}
\newcommand{\ablationFrozenBoundary}{25.82}
\newcommand{\ablationNoAugDelta}{0.35}
\newcommand{\ablationNoAugBoundary}{29.27}
\newcommand{\ablationResnetDelta}{0.07}
\newcommand{\ablationResnetBoundary}{30.70}
\newcommand{\ablationDiceDelta}{0.23}
\newcommand{\ablationDiceBoundary}{30.13}
\newcommand{\ablationLowResDelta}{5.76}
\newcommand{\ablationLowResBoundary}{26.52}
\newcommand{\repeatReferenceGap}{0.35}
\newcommand{\repeatMaxMiouGap}{0.97}
\newcommand{\repeatMaxClassGap}{6.56}
\newcommand{\repeatMaxClassName}{RO}
\newcommand{\repeatDiceDeltaFirst}{0.13}
\newcommand{\testImages}{110}
\newcommand{\deeplabTestGain}{31.52}
\newcommand{\unetTestMiou}{36.75}
\newcommand{\unetTestDice}{47.84}
\newcommand{\unetValTestGap}{0.19}
\newcommand{\unetTestIouBW}{84.82}
\newcommand{\unetTestIouHD}{33.75}
\newcommand{\unetTestIouPF}{0.05}
\newcommand{\unetTestIouWR}{33.30}
\newcommand{\unetTestIouRO}{0.00}
\newcommand{\unetTestIouRI}{54.31}
\newcommand{\unetTestIouFV}{31.64}
\newcommand{\unetTestIouSR}{56.12}
\newcommand{\deeplabCeTestMiou}{68.27}
\newcommand{\deeplabCeTestDice}{79.47}
\newcommand{\deeplabCeValTestGap}{6.22}
\newcommand{\deeplabCeTestIouBW}{89.43}
\newcommand{\deeplabCeTestIouHD}{78.96}
\newcommand{\deeplabCeTestIouPF}{24.12}
\newcommand{\deeplabCeTestIouWR}{70.32}
\newcommand{\deeplabCeTestIouRO}{71.09}
\newcommand{\deeplabCeTestIouRI}{69.57}
\newcommand{\deeplabCeTestIouFV}{70.08}
\newcommand{\deeplabCeTestIouSR}{72.58}
\newcommand{\deeplabDiceTestMiou}{69.18}
\newcommand{\deeplabDiceTestDice}{80.78}
\newcommand{\deeplabDiceValTestGap}{6.90}
\newcommand{\deeplabDiceTestIouBW}{90.04}
\newcommand{\deeplabDiceTestIouHD}{77.09}
\newcommand{\deeplabDiceTestIouPF}{34.78}
\newcommand{\deeplabDiceTestIouWR}{69.53}
\newcommand{\deeplabDiceTestIouRO}{68.87}
\newcommand{\deeplabDiceTestIouRI}{69.40}
\newcommand{\deeplabDiceTestIouFV}{70.44}
\newcommand{\deeplabDiceTestIouSR}{73.30}
\newcommand{\diceTestDeltaMiou}{0.91}
\newcommand{\diceTestDeltaBW}{0.61}
\newcommand{\diceTestDeltaHD}{1.88}
\newcommand{\diceTestDeltaPF}{10.66}
\newcommand{\diceTestDeltaWR}{0.79}
\newcommand{\diceTestDeltaRO}{2.22}
\newcommand{\diceTestDeltaRI}{0.16}
\newcommand{\diceTestDeltaFV}{0.37}
\newcommand{\diceTestDeltaSR}{0.72}
\newcommand{\testPlantsAsReef}{43.39}
\newcommand{\testRobotsAsDivers}{15.57}
\newcommand{\closeBits}{8}
\newcommand{\nearestValMedian}{15}
\newcommand{\nearestValClose}{3.5}
\newcommand{\imageMiouValMedian}{54.4}
\newcommand{\nearestTestMedian}{15}
\newcommand{\nearestTestClose}{4.5}
\newcommand{\imageMiouTestMedian}{57.7}
\newcommand{\diceDropEpoch}{26}
\newcommand{\diceDropPoints}{10.4}


\subsection{Training behavior}
\label{sec:training-behavior}
\figref{fig:unet-curves} shows the U-Net run. Training loss falls steadily
from \unetFirstTrainLoss{} in the first epoch to \unetLastTrainLoss{} in the
last, and validation loss follows it closely: it is lower than the training
loss in \unetValBelowTrain{} of the \unetEpochs{} epochs and never exceeds it
by more than \unetLossGap{}. Validation loss can be the lower of the two because
the training loss is averaged over an epoch in which the weights are still
improving, and because batch normalization uses batch statistics during
training but running statistics during evaluation. Validation mIoU stays low
during the warm-up, at \unetWarmupValMiou\% after epoch \unetWarmupEpochs{},
then rises steadily and peaks at \unetBestValMiou\% in epoch \unetBestEpoch{},
which is the selected checkpoint. The last epochs differ from this peak by only
\unetLastGap{} percentage points, and validation loss is also lowest in epoch
\unetLowestValLossEpoch{}.

The curves show no sign of overfitting. They show underfitting instead: both
losses were still falling when the cosine schedule brought the learning rate
close to zero, and the mIoU curve flattens only in these final epochs. At the
selected checkpoint the model still predicts no pixel of \unetZeroClasses{}
correctly on the validation set, so both classes have an IoU of zero; they are
the two classes with the fewest pixels (\tabref{tab:eda-classes}). A longer
schedule or a larger learning rate would probably help, but we kept the
configuration fixed in the proposal. Training took \unetTrainMinutes{} minutes,
about \unetEpochSeconds{} seconds per epoch. The only implementation issue
found for this run, a configuration changed after training, is described in
\secref{sec:training}.

\reportfigure{Images/unet_learning_curves.png}{U-Net training loss and
validation mIoU on SUIM, seed 42.}{fig:unet-curves}

\reporttodo{DeepLabV3 cross-entropy and cross-entropy plus Dice: describe both
curves, the selected epochs, overfitting or underfitting and training time, and
replace the DeepLabV3 figure, which still shows the earlier two-GPU run, with
the reported single-GPU run. State why the earlier run was replaced.}

\reportfigure{Images/deeplabv3_ce_learning_curves.png}{DeepLabV3 with
cross-entropy: training loss and validation mIoU on SUIM, seed 42.}{fig:deeplab-curves}

\reportfigure{Images/deeplabv3_dice_learning_curves.png}{DeepLabV3 with
cross-entropy plus Dice: training loss and validation mIoU on SUIM, seed 42.}{fig:dice-curves}

\subsection{Quantitative results}
\label{sec:quantitative}
\reporttodo{Report test mIoU and mean Dice for every model, using the
checkpoint selected on validation. Give per-class IoU, parameter counts,
training time, and inference cost. State that each configuration uses one seed,
or give mean and standard deviation if more seeds are run.}

\begin{table}[htbp]
\centering\small
\caption{SUIM test results with the checkpoint selected on validation. Each
configuration uses seed 42.}
\label{tab:test-results}
\begin{tabular}{lrrrr}
\toprule
Model & Loss & mIoU (\%) & Mean Dice (\%) & Parameters\\
\midrule
U-Net & CE & \pending & \pending & \pending\\
DeepLabV3 & CE & \pending & \pending & \pending\\
DeepLabV3 & CE + Dice & \pending & \pending & \pending\\
\bottomrule
\end{tabular}
\end{table}

% Generated by scripts/report/a2_results_tables.py from the U-Net training summary and the ablation results. Do not edit.
\begin{table}[htbp]
\centering\footnotesize
\setlength{\tabcolsep}{3pt}
\caption{Per-class validation IoU (\%) on the same 287 SUIM images.}
\label{tab:per-class-iou}
\begin{tabular}{lrrrrrrrr}
\toprule
Model & BW & HD & PF & WR & RO & RI & FV & SR\\
\midrule
U-Net (CE) & 80.62 & 30.83 & 0.00 & 41.16 & 0.00 & 67.94 & 33.14 & 41.85\\
DeepLabV3 (CE) & 84.19 & 70.37 & 12.00 & 69.08 & 71.66 & 74.47 & 66.84 & 50.59\\
DeepLabV3 (CE + Dice) & 83.61 & 68.49 & 12.40 & 66.92 & 74.76 & 74.04 & 67.59 & 50.38\\
\bottomrule
\end{tabular}
\end{table}


\subsection{Compute cost}
\label{sec:compute}
\tabref{tab:compute-cost} compares the cost of the models. DeepLabV3 has
\parameterRatio{} times as many parameters as U-Net but needs only \flopRatio{}
times as many operations per image, \deeplabGflops{} against \unetGflops{}
GFLOPs. U-Net spends many operations per parameter because its first and last
blocks run at the full input resolution, while DeepLabV3 runs most of its
layers at one eighth of it. The operations are counted with the PyTorch FLOP
counter for one image, with a multiply and an add counted separately. They do
not depend on the hardware, so they give the fairest comparison of the two
architectures.

Training time, latency and peak memory were recorded during the training runs.
U-Net took \unetTrainMinutes{} minutes to train and \unetLatency{} ms per image
to predict, and used at most \unetPeakMemory{}~GB of GPU memory during
training. DeepLabV3 with cross-entropy took \deeplabCeTrainMinutes{} minutes and
\deeplabCeLatency{} ms per image, and used \deeplabCePeakMemory{}~GB. The two
models were trained on different GPUs, so these times measure the hardware as
much as the models and cannot be compared directly; we report them to document
what each run costs to reproduce. Adding the Dice term changes neither the
parameters nor the operations of the network and lengthens DeepLabV3 training
by only \diceTimeOverhead\%.
\reporttodo{Name the GPU of the DeepLabV3 runs in \tabref{tab:compute-cost}.
If possible, measure the latency of both models on the same GPU.}

% Generated by scripts/report/a2_results_tables.py from the U-Net training summary and the ablation results. Do not edit.
\begin{table}[htbp]
\centering\small
\caption{Compute cost. GFLOPs are counted for one $\suimInputHeight\times\suimInputWidth$ image and do not depend on hardware. Training time, latency
and peak training memory are measured on the hardware listed; latency excludes
data loading and is measured on the validation set at batch size 4.}
\label{tab:compute-cost}
\footnotesize\setlength{\tabcolsep}{3pt}
\begin{tabular}{lrrrrrl}
\toprule
& & & Train time & Latency & Peak memory & Training\\
Model & Params (M) & GFLOPs & (min) & (ms/image) & (GB) & hardware\\
\midrule
U-Net & 7.76 & 28.3 & 14.9 & 4.8 & 0.97 & RTX 4060 Laptop\\
DeepLabV3 (CE) & 39.64 & 95.8 & 21.7 & 7.8 & 1.93 & \pending\\
DeepLabV3 (CE + Dice) & 39.64 & 95.8 & 21.9 & 7.8 & 1.93 & \pending\\
\bottomrule
\end{tabular}
\end{table}


\subsection{Qualitative results}
\label{sec:qualitative}
\reporttodo{Show test images with ground truth and predictions from each model:
correct predictions, difficult but correct cases, and failures. State the rule
used to select the examples so the selection is reproducible.}

\reportfigure{Images/a2_test_examples.png}{Test images, ground truth, and
predictions from U-Net and DeepLabV3. State the selection rule.}{fig:examples}

\subsection{Controlled experiment: loss function}
\label{sec:ablation}
\reporttodo{Report the cross-entropy against cross-entropy plus Dice result on
DeepLabV3. Compare mIoU, mean Dice, and the IoU of the rare classes RO, PF and
WR. State whether the hypothesis is supported under the decision criterion in
\secref{sec:experiments}.}

% Generated by scripts/report/a2_analysis_figures.py from comparisons.csv.
\begin{table}[htbp]
\centering\small
\caption{Controlled comparisons on the same validation split, with seed 42. Scores are percentages; $\Delta$ mIoU is the change from the reference in percentage points.}
\label{tab:ablation}
\begin{tabular}{lrrrr}
\toprule
Configuration & mIoU & Mean Dice & $\Delta$ mIoU & Boundary mIoU\\
\midrule
Reference & 62.05 & 73.91 & +0.00 & 30.07\\
Frozen backbone & 54.40 & 67.62 & -7.65 & 25.82\\
No augmentation & 61.70 & 73.75 & -0.35 & 29.27\\
ResNet101 & 62.12 & 73.82 & +0.07 & 30.70\\
CE + Dice & 62.28 & 74.42 & +0.23 & 30.13\\
Half resolution & 56.29 & 69.32 & -5.76 & 26.52\\
\bottomrule
\end{tabular}
\end{table}


\subsection{Error analysis}
\label{sec:errors}
\reporttodo{Group the errors, for example PF predicted as RI, boundary errors,
small objects missed, and background confusion. For each group give a
description, its frequency or rate from the confusion matrix, examples, a
hypothesized cause, and an improvement idea.}

\reportfigure{Images/deeplabv3_test_confusion_matrix.png}{Row-normalized
confusion matrix of the best model on the SUIM test set.}{fig:confusion}

\subsection{Comparison and discussion}
\label{sec:discussion}
\reporttodo{Explain which model is better and by how much, where the
differences appear (which classes and region sizes), and likely causes in data,
architecture, loss, or training. Discuss accuracy, speed, and parameter
trade-offs. Say which conclusions cannot yet be drawn. Listing the highest score
is not enough.}

\subsection{Limitations and conclusion}
\label{sec:conclusion}
Several limits restrict what our experiments can show. Each configuration was
trained with one seed, so differences of about one percentage point in mIoU,
and larger ones for rare classes, may come from the initialization and the
batch order rather than from the factor that was changed; the validation set
holds only \suimImagesValRO{} RO images. The hyperparameters were fixed in
advance rather than tuned, and U-Net had not converged, so the gap between the
two models partly reflects the training budget. The comparison of U-Net and
DeepLabV3 changes architecture, size and initialization together, so its result
cannot be attributed to pretraining alone. Training time and latency were
measured on different GPUs.

The data impose further limits. All metrics are computed on masks resized to
$\suimInputHeight\times\suimInputWidth$, so thin structures lose pixels and our
scores are not directly comparable with results at native resolution, and images
with other aspect ratios are stretched. SUIM covers a limited set of sites and
reuses images from EUVP, USR-248 and UFO-120, and our duplicate search misses
mirrored copies and near-identical video frames, so the test set may contain
scenes related to training scenes and the scores may not carry over to other
waters. Rejecting the masks with off-code pixels also removed some genuine
examples of rare classes. Finally, each image has one annotation and no void
label, so part of the boundary error reflects label uncertainty rather than
model mistakes.

Natural next steps are to train every configuration with several seeds and
report the mean and standard deviation, to give U-Net a longer schedule and a
small learning-rate search, to train or evaluate at a higher resolution, to add
colour and turbidity augmentation, and to make the duplicate search invariant
to mirroring.
\reporttodo{End with a short answer to both parts of the research question,
based on the test results.}
